% Финансовый прогноз на 3 года
\begin{figure}[htbp]
  \centering
  \begin{tikzpicture}
    \begin{axis}[
      width=0.9\textwidth,
      height=0.55\textwidth,
      xlabel={Месяц},
      ylabel={тыс. руб./мес.},
      xmin=0, xmax=36,
      ymin=0, ymax=450,
      grid=major,
      grid style={dashed, gray!30},
      legend style={at={(0.3,0.95)}, anchor=north, legend columns=1},
      xtick={0,6,12,18,24,30,36},
    ]
    % Выручка
    \addplot[blue, thick, mark=*, mark size=2pt] coordinates {
      (0, 0)
      (3, 15)
      (6, 45)
      (9, 75)
      (12, 120)
      (15, 150)
      (18, 180)
      (21, 210)
      (24, 240)
      (27, 270)
      (30, 300)
      (33, 330)
      (36, 360)
    };
    \addlegendentry{Выручка}
    % Расходы
    \addplot[red!70!black, thick, dashed, mark=square*, mark size=2pt] coordinates {
      (0, 200)
      (6, 200)
      (12, 200)
      (18, 200)
      (24, 200)
      (30, 200)
      (36, 200)
    };
    \addlegendentry{Расходы}
    % Точка безубыточности — вертикальная линия
    \draw[gray, densely dashed, thick] (axis cs:32,0) -- (axis cs:32,420);
    \node[font=\footnotesize, align=center, anchor=south west] at (axis cs:32.5,350)
      {Точка\\безубыточности\\(32 мес.)};
    % Точка пересечения
    \addplot[only marks, mark=o, mark size=5pt, thick, black] coordinates {(32, 200)};
    \end{axis}
  \end{tikzpicture}
  \caption{Финансовый прогноз на 3 года (сценарная модель)}
  \label{fig:financial-forecast}
\end{figure}
